% =====================================================================
%  Bolton Howes — Curriculum Vitae
%  Rebuilt from the November 2026 PDF; layout matches the original.
%
%  Editing notes
%   * Each section is a self-contained block between "% ----" rules,
%     so sections can be reordered or dropped for tenure / grant /
%     job-application versions by moving whole blocks.
%   * \me      -> bold "Bolton Howes" in author lists
%   * \stu     -> student-advisee marker (superscript asterisk)
%   * \doi{..} -> clickable DOI
%   * \pub     -> numbered publication. In-review and published papers
%     share one descending sequence; add a new paper with \pub at the
%     top of its list and all numbers update automatically (compile
%     twice). In-prep papers use plain \item and stay unnumbered.
% =====================================================================
\documentclass[10pt,letterpaper]{article}

\usepackage[T1]{fontenc}
\usepackage[utf8]{inputenc}
\usepackage{lmodern}
\usepackage[left=0.5in,right=0.5in,top=0.75in,bottom=0.75in]{geometry}
\usepackage{array}
\usepackage{longtable}
\usepackage{titlesec}
\usepackage{enumitem}
\usepackage{needspace}
\usepackage[hidelinks]{hyperref}

\pagestyle{plain}
\hyphenation{Edmonsond Moniruzzaman Duch\^atellier}
\setlength{\parindent}{0pt}
\renewcommand{\arraystretch}{1.25}

% ---------- Section headings: large sans-serif bold with full rule ----
\titleformat{\section}{\Large\sffamily\bfseries}{}{0em}{}[\vspace{3pt}\titlerule]
\titlespacing*{\section}{0pt}{14pt}{6pt}
\newcommand{\cvsection}[2][]{\Needspace{6\baselineskip}\section*{#2#1}}
\newcommand{\cvnote}[1]{\hspace{0.4em}{\normalfont\rmfamily\footnotesize #1}}
\newcommand{\cvsub}[1]{\par\Needspace{4\baselineskip}\vspace{4pt}\textit{#1}\par\vspace{2pt}}

% ---------- Shorthands ------------------------------------------------
\newcommand{\me}{\textbf{Bolton Howes}}
\newcommand{\stu}{\textsuperscript{*}}
\newcommand{\doi}[1]{\href{https://doi.org/#1}{doi:#1}}
\newcommand{\twoyears}[2]{{\small #1 {\footnotesize\&} #2}}

% ---------- Tables: date column | rule | content ----------------------
\setlength{\LTleft}{0pt}
\setlength{\LTright}{0pt}
\setlength{\LTpre}{2pt}
\setlength{\LTpost}{0pt}
\setcounter{LTchunksize}{1}  % let long tables break at any row
\newlength{\datecol}   \setlength{\datecol}{1.75cm}
\newlength{\namecol}   \setlength{\namecol}{3.5cm}
\newlength{\placecol}  \setlength{\placecol}{2.75cm}
\newlength{\weekcol}   \setlength{\weekcol}{1.5cm}
\newlength{\restA}
\setlength{\restA}{\dimexpr\textwidth-\datecol-4\tabcolsep-\arrayrulewidth\relax}
\newlength{\restB}
\setlength{\restB}{\dimexpr\textwidth-\datecol-\namecol-6\tabcolsep-\arrayrulewidth\relax}
\newlength{\restC}
\setlength{\restC}{\dimexpr\textwidth-\datecol-\namecol-\weekcol-8\tabcolsep-\arrayrulewidth\relax}
\newlength{\restD}
\setlength{\restD}{\dimexpr\textwidth-\datecol-\placecol-\weekcol-8\tabcolsep-\arrayrulewidth\relax}

% date | text
\newenvironment{cvtable}{\longtable{p{\datecol}|p{\restA}}}{\endlongtable}
% date | name | text
\newenvironment{cvtableB}{\longtable{p{\datecol}|p{\namecol}p{\restB}}}{\endlongtable}
% date | name | weeks | place
\newenvironment{cvtableC}{\longtable{p{\datecol}|p{\namecol}p{\weekcol}p{\restC}}}{\endlongtable}
% date | place | weeks | description
\newenvironment{cvtableD}{\longtable{p{\datecol}|p{\placecol}p{\weekcol}p{\restD}}}{\endlongtable}

% ---------- Publication lists -----------------------------------------
\newenvironment{publist}%
  {\begin{list}{}{\setlength{\leftmargin}{2.5em}\setlength{\labelwidth}{2em}%
     \setlength{\labelsep}{0.5em}%
     \setlength{\itemsep}{6pt}\setlength{\topsep}{4pt}\setlength{\parsep}{0pt}}}%
  {\end{list}}
% Descending numbering across lists: the total is saved to the .aux file
\newcounter{pubdown}
\newcounter{pubcount}
\makeatletter
\providecommand{\totalpubs}{0}
\AtBeginDocument{\setcounter{pubdown}{\totalpubs}}
\AtEndDocument{\immediate\write\@auxout{\gdef\string\totalpubs{\the\value{pubcount}}}}
\makeatother
\newcommand{\pub}{\item[\arabic{pubdown}.]\addtocounter{pubdown}{-1}\stepcounter{pubcount}}

\begin{document}

% =====================================================================
%  HEADER
% =====================================================================
{\centering
  {\LARGE\bfseries Bolton Howes}\par\vspace{7pt}
  \textit{Assistant Professor \& Director of the Wheaton College Field Station (Black Hills, SD)}\par\vspace{2pt}
  Wheaton College\par\vspace{6pt}
  \href{mailto:bolton.howes@wheaton.edu}{bolton.howes@wheaton.edu} $|$
  \href{https://boltonhowes22.github.io}{boltonhowes22.github.io}\par}

% ---------------------------------------------------------------------
\cvsection{Professional Experience}
\begin{cvtable}
2025--      & \textbf{Wheaton College}\newline
              Assistant Professor \& Director of the Wheaton College Field Station (Black Hills, SD) \\
2024--2025  & \textbf{University of Washington \& Western Washington University}\newline
              NSF Postdoctoral Fellow\newline
              Advisors: Akshay Mehra and Brady Foreman \\
2023--2024  & \textbf{University of Miami (Rosenstiel School of Marine, Atmospheric, and Earth Science)}\newline
              Postdoctoral Researcher, Marine Geology\newline
              Advisor: Sam Purkis \\
\end{cvtable}

% ---------------------------------------------------------------------
\cvsection{Education}
\begin{cvtable}
2017--2023  & \textbf{Princeton University}, Princeton, NJ\newline
              Ph.D., Geosciences\newline
              Dissertation: Developing physical paleoenvironmental proxies: from ooids to extinction\newline
              Advisor: Adam Maloof \\
2015--2017  & \textbf{University of Georgia}, Athens, GA\newline
              M.Sc., Geology\newline
              Thesis: Sequence stratigraphic expression of flexural subsidence: Middle Jurassic Twin Creek Limestone, Wyoming, U.S.A.\newline
              Advisor: Steven Holland \\
2011--2015  & \textbf{Macalester College}, St. Paul, MN\newline
              B.A., Geology\newline
              Advisor: Raymond Rogers \\
\end{cvtable}

% ---------------------------------------------------------------------
\cvsection{Awards and Honors}
\begin{cvtable}
2024        & UW eScience Postdoctoral Fellowship [\$4,000] \\
2023        & NSF EAR Postdoctoral Fellowship [\$180,000] \\
2017--2022  & Department of Geosciences Graduate Research Fund [3x\$5,000] \\
2017        & Outstanding M.S. Student of the Year, Department of Geology, University of Georgia \\
2015--2017  & Graduate School Assistantship, University of Georgia (university-wide competition; department nomination, faculty-panel selection) \\
2016        & Graduate Research Grant, Geological Society of America \\
2016        & Miriam Watts-Wheeler Award, University of Georgia [\$1,600] \\
2015        & Henry Lepp Award for Dedication to Scientific Research, Macalester College \\
\end{cvtable}

% ---------------------------------------------------------------------
\cvsection[\cvnote{(\stu\ denotes student advisee)}]{Publications}

\cvsub{In review}
\begin{publist}
  \pub Filip Novak\stu, \me, Henry Yuan, Brady Foreman, Akshay Mehra. From outcrop to point cloud: A three-dimensional technique for the remote analysis of terrestrial and extraterrestrial sedimentary structures. In review at \textit{GSA Today}.
  \pub Anthony Semeraro, Benjamin Cardenas, \me, Brady Foreman. Reconstructing river planform morphology in alluvial stratigraphy using paleocurrents. In review at \textit{The Depositional Record}.
\end{publist}

\cvsub{Published}
\begin{publist}
  \pub Morgan Chakraborty, Hannah Shernisky, Clément Pollier, Amanda Oehlert, Arash Sharifi, Ali Pourmand, \me, Benjamin Minch, Mohammad Moniruzzaman, Larissa Frühe, Mattie Rodrigue, Sam Purkis. Brine pool microbes enrich metalliferous sediments in salt giant basins. 2026. \textit{AGU Advances}. \doi{10.1029/2026AV002570}
  \pub Kevin Vélez-Rosado, Olga Zalles-Grebetskaya, Jeffrey Wilson Mantilla, Blair Schoene, Adam Maloof, \me. New material of \textit{Dolichochampsa minima} (Archosauria: Crocodylia) from the Cretaceous--Palaeogene El Molino Formation of Bolivia sheds light on the early evolution of Gavialinae. 2025. \textit{Journal of Systematic Palaeontology}. \doi{10.1080/14772019.2025.2496524}
  \pub Morgan Chakraborty, Arash Sharifi, Francesca Benzoni, François Tissot, Ali Pourmand, Marco Taviani, \me, Peter Swart, Chaojin Lu, Mattie Rodrigue, Sam Purkis. Deep-water corals indicate the Red Sea survived the last glacial lowstand. 2025. \textit{PNAS}. \doi{10.1073/pnas.2415559122}
  \pub Sam Purkis, Steven Ward, \me, Jake Longenecker, Morgan Chakraborty, Akos Kalman, Amy Clement, Arash Sharifi, Francesca Benzoni, Christopher Clarke, Mattie Rodrigue. A 1600-year record of extreme rainfall in northern Arabia. 2025. \textit{Science Advances}. \doi{10.1126/sciadv.adq3173}
  \pub \me, Akshay Mehra, Emily Geyman, Julia Wilcots, Ryan Manzuk, Curtis Deutsch, Adam Maloof. The where, when, and how of ooid formation: what ooids tell us about ancient seawater chemistry. 2024. \textit{Earth and Planetary Science Letters}. \doi{10.1016/j.epsl.2024.118697}
  \pub Gaëlle Duchâtellier, Amanda Oehlert, Hannah Shernisky, Clément Pollier, Peter Swart, \me, Sam Purkis. Sedimentary porewaters record regional tectonic and climate events that perturbed a deep-sea brine pool in the Gulf of Aqaba, Red Sea. 2024. \textit{Science of the Total Environment}. \doi{10.1016/j.scitotenv.2023.168804}
  \pub Akshay Mehra, \me, Ryan Manzuk, Alex Spatzier, Bradley Samuels, Adam Maloof. A novel technique for producing three-dimensional data using serial sectioning and semi-automatic image classification. 2022. \textit{Microscopy and Microanalysis}. \doi{10.1017/S1431927622012442}
  \pub Emily Geyman, Ziman Wu, Matthew Nadeau, Stacey Edmonsond, Andrew Turner, Sam Purkis, \me, Blake Dyer, Anne-Sofie Ahm, Nan Yao, Curtis Deutsch, John Higgins, Daniel Stolper, Adam Maloof. The origin of carbonate mud and implications for global climate. 2022. \textit{PNAS}. \doi{10.1073/pnas.2210617119}
  \pub \me, Akshay Mehra, Adam Maloof. Three-dimensional morphometry of ooids in oolites: a new tool for more accurate and precise paleoenvironmental interpretation. 2021. \textit{Journal of Geophysical Research: Earth Surface}. \doi{10.1029/2020JF005601}
  \pub Paul Hoffman, Galen Halverson, Daniel Schrag, John Higgins, \ldots, \me, et al. Snowballs in Africa: sectioning a long-lived Neoproterozoic carbonate platform and its bathyal foreslope (NW Namibia). 2021. \textit{Earth-Science Reviews}. \doi{10.1016/j.earscirev.2021.103616}
\end{publist}

\cvsub{In prep.}
\begin{publist}
  \item Filip Novak\stu, \me, Brady Foreman, Akshay Mehra. Minimal hydrologic changes spanning the Paleocene-Eocene Thermal Maximum: Biased proxy or buffered signal?
  \item \me, Morgan Chakraborty, Sam Purkis. The afterlife of lowstand sediments.
  \item \me, Delaney Todd, Katherine Sullivan, Akshay Mehra, Brady Foreman. Increased river mobility and floodplain erosion during the Paleocene-Eocene global warming event.
\end{publist}

% ---------------------------------------------------------------------
\cvsection[\cvnote{(\stu\ denotes student advisee)}]{Conference Abstracts}
\begin{cvtable}
2026 & Filip Novak\stu, \me, Brady Foreman, Akshay Mehra. Cross-sets as a Paleohydrologic Proxy: Insights from the Paleocene-Eocene Thermal Maximum of the Bighorn Basin, Wyoming. \textit{GSA} (2026). \\
2025 & Filip Novak\stu, \me, Brady Foreman, Akshay Mehra. From Outcrop to Point Cloud: Using Three-Dimensional Techniques to Study Fluvial Dynamics During the Paleocene-Eocene in Wyoming. \textit{GSA} (2025). \\
2025 & \me, Anthony Semeraro, Akshay Mehra, Brady Foreman. Reconstructing river planform morphology. \textit{Geobiology} (2025). \\
2022 & \me, Ryan Manzuk, Akshay Mehra, Adam Maloof. Tales of transport and chemistry from three-dimensional growth histories of ooids. \textit{AGU Fall Meeting} (Chicago, Illinois). \\
2022 & Ryan Manzuk, \me, Emily Geyman, Indu Panigrahi, Dave Singh, Adam Maloof. Multispectral petrography image analysis as a method for paleoenvironmental and paleoecological reconstruction; applications to Earth's first reefs. \textit{AGU Fall Meeting} (Chicago, Illinois). \\
2022 & Cedric Hagen, John Higgins, Adam Maloof, Stacey Edmonsond, Ryan Manzuk, \me, Pingping Zhao. Machine learning and computational image analysis leveraged to probe relationships between carbonate isotopic composition, diagenetic alteration, and sedimentary facies isotopic variability. \textit{AGU Fall Meeting} (Chicago, Illinois). \\
2022 & Adam Maloof, Ryan Manzuk, Emily Geyman, Akshay Mehra, Jaap Kaandorp, Mark Webster, Stacey Edmonsond, \me, and Cedric Hagen. From modern analogs to three dimensions: Lessons learned for interpreting the stratigraphic record of the Proterozoic--Phanerozoic transition. \textit{GSA Annual Meeting} (Denver, Colorado). \\
2018 & \me, Akshay Mehra, and Adam Maloof. Three-dimensional reconstructions of Holocene and Neoproterozoic oolites to measure porosity, permeability, and volume-shape evolution of ooids. \textit{GSA Annual Meeting} (Indianapolis, Indiana). \\
2014 & \me\ and Raymond Rogers. Revisiting a regionally significant terrestrial bounding surface in the Upper Cretaceous (Campanian) Two Medicine Formation, northwestern Montana. \textit{GSA Annual Meeting} (Vancouver, British Columbia). \\
\end{cvtable}

% ---------------------------------------------------------------------
\cvsection{Invited Talks}
\begin{cvtable}
2026 & University of Chicago \\
2026 & University of Washington College of the Environment, Pathways to a Professor Panel \\
2025 & Purdue University \\
2025 & University of Washington, PaleoLunch \\
2025 & Macalester College, Sed/Strat Seminar \\
2025 & Western Washington University, Department Seminar \\
2024 & Wheaton College, Department Seminar \\
2020 & Macalester College, Graduate School Open House \\
2018 & Macalester College, Sed/Strat Seminar \\
\end{cvtable}

% ---------------------------------------------------------------------
\cvsection{Fieldwork Experience (84 Weeks)}
\begin{cvtableD}
2026                   & Bolivia            & 2 weeks  & K--Pg paleontology \\
2024                   & Indonesia          & 2 weeks  & Mapping seafloor ecology and natural hazards \\
\twoyears{2024}{2025}  & NW Wyoming         & 5 weeks  & The fluvial response to the PETM \\
2023                   & Red Sea            & 3 weeks  & Geomorphology and stratigraphy to evaluate natural hazard risk \\
\twoyears{2019}{2022}  & Bolivia            & 12 weeks & The K--Pg Boundary in lacustrine and fluvial settings \\
2021                   & southern Utah      & 2 weeks  & Cyclicity in the Late Paleozoic Ice Age \\
2019                   & Calabria, Italy    & 1 week   & Measuring climate change in olive orchards \\
2019                   & Yukon, Canada      & 5 weeks  & Archaeocyathid reefs and the Cambrian Explosion \\
2018                   & Utah               & 1 week   & Measuring climate change (GEOL 201 course fieldwork) \\
2018                   & Spain and Italy    & 6 weeks  & The K--Pg Boundary and slope sediments \\
2018                   & southern Namibia   & 4 weeks  & Neoproterozoic and Paleozoic glacial deposits \\
2018                   & Western Australia  & 1 week   & Modern and recent carbonate platforms \\
2018                   & Nevada/Arizona     & 3 weeks  & Stratigraphy and geochemistry of the Late Paleozoic Ice Age \\
2017                   & northern Namibia   & 6 weeks  & Field assistant to Paul Hoffman. Mapping Snowball Earth \\
2016                   & western Wyoming    & 6 weeks  & Sequence stratigraphy of the Twin Creek Limestone \\
2014                   & NW Montana         & 2 weeks  & Sequence stratigraphy and stratigraphic paleobiology, Two Medicine Fm. \\
2014                   & Patagonia, Chile   & 12 weeks & Restoration ecology and conservation biology in the Aysén Region \\
\twoyears{2013}{2014}  & NW Wyoming         & 5 weeks  & Terrestrial stratigraphy and paleobiology of the Cloverly Formation \\
2014                   & Tall Dhiban,\newline Jordan & 6 weeks & Archeological excavation: Early Bronze Age through the 19th century \\
\end{cvtableD}

% ---------------------------------------------------------------------
\cvsection{Graduate Committees}
\begin{cvtableB}
2026 & Katherine Izzo & M.Sc. Committee Member, Western Washington University \\
\end{cvtableB}

% ---------------------------------------------------------------------
\cvsection{Undergraduates Advised}
\begin{cvtableB}
2026       & Carson Frear          & Bolivian K--Pg \\
2024--2025 & Filip Novak           & Photogrammetry and fluvial sedimentology \\
2021       & Anubhav Agarwal       & Transfer learning for geological image segmentation \\
2021       & Hugh Shields          & Transfer learning for geological image segmentation \\
2021       & Anna Krokhine         & Numerical models of ooid growth \\
2020       & Emilio Cano-Renteria  & Geological image segmentation and ooid growth \\
2019       & Sarah Brown           & K--Pg boundary stratigraphy of Bolivia \\
2018       & Celia Aranda-Reina    & Climate variability during the K--Pg in Spain \\
\end{cvtableB}

% ---------------------------------------------------------------------
\cvsection{Undergraduates Involved in Fieldwork}
\begin{cvtableC}
2026 & Carson Frear        & 2 weeks & Bolivia \\
2025 & Filip Novak         & 2 weeks & Wyoming \\
2019 & 10 Undergraduates   & 1 week  & Italy \\
2019 & Sarah Brown         & 5 weeks & Bolivia \\
2019 & Galen Cadley        & 5 weeks & Yukon, Canada \\
2018 & 8 Undergraduates    & 1 week  & Utah \\
2018 & Celia Aranda-Reina  & 4 weeks & Spain \\
2016 & Tony Moraes         & 4 weeks & western Wyoming \\
\end{cvtableC}

% ---------------------------------------------------------------------
\Needspace{10\baselineskip}% keep heading with its first subsection
\cvsection{Teaching}

\cvsub{Instructor of Record}
\begin{cvtable}
Spr. 2027 & \textbf{GEOL 212:} The Dynamic Earth and Environment, Wheaton College \\
Spr. 2026 & \textbf{GEOL 321:} Earth History and Stratigraphy, Wheaton College \\
\end{cvtable}

\cvsub{Laboratory Instructor}
\begin{cvtable}
Spr. 2026 & \textbf{GEOL 212L:} The Dynamic Earth and Environment (lab section), Wheaton College \\
Fall 2025 & \textbf{GEOL 212L:} The Dynamic Earth and Environment (lab section), Wheaton College \\
\end{cvtable}

\cvsub{Teaching Assistant}
\begin{cvtable}
Spr. 2021            & \textbf{ENV 354-GEO 368} Climate and Weather: Order in the Chaos, Princeton University \\
Fall 2019\textsuperscript{\textbf{*}} & \textbf{FRS 161:} Earth: Crops, Culture, and Climate (in Italy), Princeton University \\
Fall 2018\textsuperscript{\textbf{*}} & \textbf{GEOL 201:} Measuring Climate Change: Data Analysis and Science Writing, Princeton University \\
2016--2017           & Geology Tutor, University of Georgia Student-Athlete Academic Center \\
Spr. 2015            & \textbf{GEOL 165:} History and Evolution of the Earth, Macalester College \\
Fall 2014            & \textbf{GEOL 300:} Paleobiology, Macalester College \\
\end{cvtable}

\vspace{4pt}
\textit{\textsuperscript{\textbf{*}} Indicates courses that included a fieldwork component consisting of a week-long trip (Utah and Italy) to make observations for research projects. My teaching duties included preparing supplies and logistics, and helping students with both data collection---especially with drones, differential GPS, and surveying equipment---and data analysis. The data collected during fieldwork was part of a semester-long project.}

% ---------------------------------------------------------------------
\cvsection{Leadership and Service}
\begin{cvtable}
2023       & Fellow, Scholars Institute Fellows Program (mentoring first-generation students), Princeton University \\
2018--2021 & Residential Graduate Student, Princeton University \\
2016--2017 & Student Volunteer, Geological Society of America \\
2015       & Leadership Committee of Macalester College Fellowship of Christian Athletes \\
2012--2013 & Resident Assistant, Macalester College \\
2011--2014 & Captain of Macalester College Football Team \\
\end{cvtable}

% ---------------------------------------------------------------------
\cvsection{Professional Activities}
\begin{cvtable}
Convener & 2019 GSA Annual Meeting technical session: Hello (ancient) world!: Exploring the Neoproterozoic to Cambrian interval by quantitatively probing the rock record \\
Reviewer & NSF EAR Geoinformatics, Geophysical Research Letters \\
\end{cvtable}

% ---------------------------------------------------------------------
\cvsection{Public Engagement}
\begin{cvtable}
2025 & \textit{Tiny Living Beings} podcast (American Society for Microbiology): ``4.5 Billion Years of Earth History in One Hour'' \\
\end{cvtable}

\end{document}
